\documentclass[12pt,a4paper]{article}
\usepackage{graphicx}
\usepackage{amsfonts,amsmath}
\usepackage{listings}
\usepackage[T1]{fontenc}
\usepackage{float}

\title{\textbf{Sistemas Operativos: Práctica 3}}
\author{
    Shaofan Xu \\
    \and 
    Javier Maria Santa
    }
\date{\today}


\begin{document}
\maketitle

\clearpage
\section{Introducción}
En esta práctica el objetivo es ampliar la simulación de minería de blockchain introduciendo mecanismos avanzados de comunicación entre procesos (IPC) y memoria compartida. En concreto, se centraliza la salida de los mineros a través de un sistema de monitorización basado en colas de mensajes y un esquema productor-consumidor.

\section{Requisitos cumplidos y no cumplidos}
\subsection{Requisitos cumplidos}
Se han implementado correctamente los siguientes objetivos principales de la práctica, correspondientes a los apartados a y b:
\begin{itemize}
    \item \textbf{Sistema de Monitorización (Apartado a):} Se ha desarrollado el ejecutable \texttt{monitor}, el cual exige estar arrancado antes que cualquier minero. Este consta de dos procesos que gestionan la salida unificada:
    \begin{itemize}
        \item \textbf{Comprobador (Padre):} Recibe los bloques ganadores a través de una cola de mensajes POSIX (\texttt{/mq\_monitor}). Antes de procesarlos, verifica su validez comprobando que $\mathrm{pow\_hash(solucion)} == \mathrm{target}$. Si es válido, actúa como \textit{Productor} e introduce el bloque en memoria compartida.
        \item \textbf{Monitor (Hijo):} Actúa como \textit{Consumidor}, leyendo los bloques de la memoria compartida e imprimiéndolos por pantalla respetando los \textit{lags} de tiempo exigidos en la invocación.
    \end{itemize}
    \item \textbf{Esquema Productor-Consumidor IPC:} La comunicación interna del Monitor se ha sincronizado empleando un segmento de memoria compartida (\texttt{shm\_open} y \texttt{mmap}) que alberga un búfer circular. El acceso concurrente a dicho búfer está protegido mediante tres semáforos POSIX con nombre(que es sustituido por semaforos sin nombres): \texttt{mutex}, \texttt{empty} y \texttt{full}.
    \item \textbf{Adaptación del Minero (Apartado b):} El proceso Minero envía un mensaje con su pid, ronda, target y solución por la cola al proclamarse ganador de una votación. Asimismo, cuando el tiempo del minero expira y detecta que es el último en salir, envía un bloque especial de terminación (\texttt{\{-1, -1\}}) por la cola.
    \item \textbf{Limpieza controlada de recursos:} Al recibir el bloque de finalización del último minero, el Comprobador intercepta el mensaje y se encarga de realizar las llamadas \texttt{mq\_unlink}, \texttt{shm\_unlink} y \texttt{sem\_unlink}, cerrando a su vez a su proceso hijo. Esto garantiza que no existan fugas de memoria en las estructuras del sistema operativo, como se ha comprobado en los test con Valgrind.
\end{itemize}

\subsection{Requisitos no cumplidos}
Todos los requisitos obligatorios especificados en el enunciado han sido implementados 
y funcionan correctamente.

\section{Pruebas: Funcionamiento y Rendimiento}
\subsection{Entorno de ejecución}
Las pruebas y mediciones se han llevado a cabo bajo las siguientes condiciones de
hardware y software para garantizar la consistencia de los resultados:
\begin{itemize}
    \item Sistema Operativo: WSL2 (Ubuntu) sobre Windows 11.
    \item Hardware: 16 GB de memoria RAM.
    \item Herramienta de medición: Comando time de la shell de Linux.
\end{itemize}

\subsection{Pruebas de Funcionamiento}
Se han ejecutado diversas pruebas para asegurar el correcto funcionamiento del programa.

\subsubsection{Casos normales}
Se ejecutó el programa con parámetros correspondientes (ej.target inicial 100, 10 rondas y 4 hilos):
\begin{lstlisting}[language=bash]
./miner 100 10 4
\end{lstlisting}
El programa generó correctamente el archivo \texttt{.log} con la estructura especificada(Id, ronda, target, solución, etc ), alternando la validación (accepted/rejected). Al finalizar, se imprime en la terminal exactamente los mensajes requeridos:
\begin{lstlisting}[language=bash]

\subsubsection{Casos de errores}
Se probaron situaciones anómalas como:
\begin{itemize}
    \item \textbf{Falta de argumentos:} Al ejecutar \texttt{./miner} sin parámetros, el programa detecta el error, muestra el uso correcto \textit{``Usage ./miner < TARGET\_INI > < ROUNDS > < N\_THREADS >''} e imprime \textit{``Miner exited unexpectedly''}.
\end{itemize}

\subsection{Pruebas de Rendimiento}
Para evaluar la eficiencia de la paralelización implementada, se han realizado pruebas de rendimiento variando el número de hilos de ejecución. Estas
pruebas permiten analizar cómo el sistema escala al dividir el espacio de búsqueda entre múltiples unidades de procesamiento.
\subsubsection{Pruebas realizadas}

\subsubsection{Resultados Obtenidos}

\subsubsection{Análisis y Razonamiento}
\section{Conclusión y Respuesta a las cuestiones del enunciado}
\textbf{Pregunta d) En base al funcionamiento del sistema anterior:}
\begin{itemize}
    \item \textbf{¿Es necesario un sistema tipo productor-consumidor para garantizar que el sistema funciona correctamente si el retraso de Comprobador es mucho mayor que el de Minero? ¿Por qué?}
    
    Sí, es absolutamente necesario. Si el \texttt{Comprobador} tiene un retraso (lag) mucho mayor que el de los \texttt{Mineros}, se convierte en el cuello de botella del sistema.
    Los mineros generarán bloques ganadores más rápido de lo que el \texttt{Comprobador} puede verificarlos. El modelo productor-consumidor (apoyado por la cola de mensajes y el búfer circular) actúa como un amortiguador (buffer).
    Permite que los mineros sigan depositando bloques en la cola sin quedarse bloqueados de inmediato, y a su vez, permite que el \texttt{Comprobador} deposite los bloques verificados en la memoria compartida de forma asíncrona respecto al \texttt{Monitor}.
    Sin este sistema, la sincronización estricta obligaría a los mineros a detener su ejecución matemática hasta que el \texttt{Comprobador} y el \texttt{Monitor} terminasen sus lentas tareas de I/O y validación.

    \item \textbf{¿Y si el retraso de Comprobador es muy inferior al de Minero? ¿Por qué?}
    
    Sí, sigue siendo conceptualmente necesario y una buena práctica de diseño de Sistemas Operativos, aunque el riesgo de colapso sea menor. Si el \texttt{Comprobador} es muy rápido, 
    el búfer circular estará vacío la mayor parte del tiempo, ya que consumirá los bloques casi al instante. Sin embargo, el paradigma productor-consumidor es vital porque \textbf{desacopla} los procesos. 
    Garantiza la exclusión mutua (evitando condiciones de carrera al acceder a la memoria compartida) y proporciona sincronización condicional (bloqueando al \texttt{Monitor} limpiamente cuando no hay bloques, en lugar de hacer espera activa consumiendo CPU). 

    \item \textbf{¿Se simplificaría la estructura global del sistema si entre Comprobador y Monitor hubiera también una cola de mensajes? ¿Por qué?}
    
    Sí, la estructura global y el código del programa se simplificarían drásticamente. Las colas de mensajes POSIX (\texttt{mq\_open}, \texttt{mq\_send}, \texttt{mq\_receive}) ya implementan internamente el patrón productor-consumidor a nivel del kernel del Sistema Operativo. 
    El kernel gestiona automáticamente la memoria del búfer interno, el control de concurrencia (mutex) y los bloqueos por cola vacía o llena. Al usar memoria compartida pura (\texttt{shm\_open}), como hemos hecho en esta práctica, el programador está obligado a gestionar manualmente
     los punteros del búfer circular (\texttt{in}, \texttt{out}) y a coordinar explícitamente los semáforos (\texttt{mutex}, \texttt{empty}, \texttt{full}), lo que incrementa la complejidad del código y el riesgo de interbloqueos (deadlocks).
\end{itemize}
\end{document}